% GENERATED FILE. Edit canonical lessons, then run npm run generate:books.
% canonical-source-hash: fnv1a64:2fdcd06f976ebe85
% canonical-lessons: BN-C381-bebostha, BN-C381-karjokor, BN-C381-bikol, BN-C381-dokkhota, BN-C381-onushilon-kora

\chapter{Making It Work}
\label{ch:bn-making-it-work}
\hlchaptermodality{381}

\begin{chapteropening}
\textbf{By the end of this chapter:} \emph{I can say what arrangements work and what breaks down.}

Complete the last lesson of chapter 381: I can say what arrangements work and what breaks down.
\end{chapteropening}

\section[\texorpdfstring{\bn{ব্যবস্থা} (bæbôsthā)}{ব্যবস্থা (bæbôsthā)}]{\bn{ব্যবস্থা} (bæbôsthā) — an arrangement; a system}

\label{lesson:BN-C381-bebostha}



\begin{quote}
Before the new one: say the Bengali for qualified, then the Bengali for to get used to.
\end{quote}

\subsection*{You'll want to know: \bn{ব্যবস্থা}}
\textbf{\bn{ব্যবস্থা}} — \emph{bæbôsthā} — \textquotedblleft{}an arrangement; a system\textquotedblright{}.

\begin{quote}
\bn{তুমি} \bn{এলে} \bn{থাকার} \bn{ব্যবস্থা} \bn{আমি} \bn{করব।}
\end{quote}

\begin{quote}
\emph{tumi ele thākār bæbôsthā āmi kôrbô.}
\end{quote}

\begin{quote}
If you come, I'll arrange a place for you to stay.
\end{quote}

\begin{cousinweb}
\textbf{\bn{ব্যবস্থা}} (\emph{bæbôsthā}) is Sanskrit \emph{vyavasthā}, an arrangement or settled order (Monier-Williams). \bn{ব্যবস্থা} \bn{করা} is to arrange something.
\end{cousinweb}

\begin{grammarlens}[title={What you've built}]
One more word for saying what happens if.
\end{grammarlens}

\subsection*{Your turn}
\begin{itemize}
\raggedright
  \item \emph{Say it:} \emph{bæbôsthā}
  \item \emph{Say it:} \emph{bæbôsthā}, once more
  \item \emph{Say it:} \emph{tumi ele thākār bæbôsthā āmi kôrbô.} — If you come, I'll arrange a place for you to stay.
  \item \emph{Recall:} say the Bengali for qualified, then the Bengali for to get used to, then say \emph{bæbôsthā} again
\end{itemize}

\begin{tcolorbox}[breakable,colback=teal!4,colframe=teal!35!black,title={Before you move on}]
What does \bn{ব্যবস্থা} mean? (\textquotedblleft{}An arrangement; a system\textquotedblright{}.) Say it once more.
\end{tcolorbox}
\section[\texorpdfstring{\bn{কার্যকর} (kārjôkôr)}{কার্যকর (kārjôkôr)}]{\bn{কার্যকর} (kārjôkôr) — effective}

\label{lesson:BN-C381-karjokor}



\begin{quote}
Before the new one: say the Bengali for to get used to, then the Bengali for an arrangement.
\end{quote}

\subsection*{You'll want to know: \bn{কার্যকর}}
\textbf{\bn{কার্যকর}} — \emph{kārjôkôr} — \textquotedblleft{}effective\textquotedblright{}.

\begin{quote}
\bn{এই} \bn{পদ্ধতি} \bn{কার্যকর} \bn{হলে} \bn{আমরা} \bn{সবাই} \bn{এটা} \bn{ব্যবহার} \bn{করব।}
\end{quote}

\begin{quote}
\emph{ei pôddhôti kārjôkôr hôle āmrā sôbāi eṭā bæbôhār kôrbô.}
\end{quote}

\begin{quote}
If this method works, we'll all use it.
\end{quote}

\begin{cousinweb}
\textbf{\bn{কার্যকর}} (\emph{kārjôkôr}) is Sanskrit \emph{kāryakara}, 'work-doing', efficacious (Monier-Williams): a method or a law that does its job.
\end{cousinweb}

\begin{grammarlens}[title={What you've built}]
One more word for saying what happens if.
\end{grammarlens}

\subsection*{Your turn}
\begin{itemize}
\raggedright
  \item \emph{Say it:} \emph{kārjôkôr}
  \item \emph{Say it:} \emph{kārjôkôr}, once more
  \item \emph{Say it:} \emph{ei pôddhôti kārjôkôr hôle āmrā sôbāi eṭā bæbôhār kôrbô.} — If this method works, we'll all use it.
  \item \emph{Recall:} say the Bengali for to get used to, then the Bengali for an arrangement, then say \emph{kārjôkôr} again
\end{itemize}

\begin{tcolorbox}[breakable,colback=teal!4,colframe=teal!35!black,title={Before you move on}]
What does \bn{কার্যকর} mean? (\textquotedblleft{}Effective\textquotedblright{}.) Say it once more.
\end{tcolorbox}
\section[\texorpdfstring{\bn{বিকল} (bikôl)}{বিকল (bikôl)}]{\bn{বিকল} (bikôl) — out of order; broken down}

\label{lesson:BN-C381-bikol}



\begin{quote}
Before the new one: say the Bengali for an arrangement, then the Bengali for effective.
\end{quote}

\subsection*{You'll want to know: \bn{বিকল}}
\textbf{\bn{বিকল}} — \emph{bikôl} — \textquotedblleft{}out of order; broken down\textquotedblright{}.

\begin{quote}
\bn{গাড়ি} \bn{বিকল} \bn{হলে} \bn{আমাকে} \bn{ফোন} \bn{করো।}
\end{quote}

\begin{quote}
\emph{gāṛi bikôl hôle āmāke phon kôro.}
\end{quote}

\begin{quote}
Call me if the car breaks down.
\end{quote}

\begin{cousinweb}
\textbf{\bn{বিকল}} (\emph{bikôl}) is Sanskrit \emph{vikala}, impaired or defective (Monier-Williams). It is the word for a machine or a body part that has stopped working.
\end{cousinweb}

\begin{grammarlens}[title={What you've built}]
One more word for saying what happens if.
\end{grammarlens}

\subsection*{Your turn}
\begin{itemize}
\raggedright
  \item \emph{Say it:} \emph{bikôl}
  \item \emph{Say it:} \emph{bikôl}, once more
  \item \emph{Say it:} \emph{gāṛi bikôl hôle āmāke phon kôro.} — Call me if the car breaks down.
  \item \emph{Recall:} say the Bengali for an arrangement, then the Bengali for effective, then say \emph{bikôl} again
\end{itemize}

\begin{tcolorbox}[breakable,colback=teal!4,colframe=teal!35!black,title={Before you move on}]
What does \bn{বিকল} mean? (\textquotedblleft{}Out of order; broken down\textquotedblright{}.) Say it once more.
\end{tcolorbox}
\section[\texorpdfstring{\bn{দক্ষতা} (dôkkhôtā)}{দক্ষতা (dôkkhôtā)}]{\bn{দক্ষতা} (dôkkhôtā) — skill; competence}

\label{lesson:BN-C381-dokkhota}



\begin{quote}
Before the new one: say the Bengali for effective, then the Bengali for out of order.
\end{quote}

\subsection*{You'll want to know: \bn{দক্ষতা}}
\textbf{\bn{দক্ষতা}} — \emph{dôkkhôtā} — \textquotedblleft{}skill; competence\textquotedblright{}.

\begin{quote}
\bn{দক্ষতা} \bn{থাকলে} \bn{চাকরি} \bn{পাওয়া} \bn{সহজ।}
\end{quote}

\begin{quote}
\emph{dôkkhôtā thākle chākri pāoyā sôhôj.}
\end{quote}

\begin{quote}
If you have the skills, it's easy to get a job.
\end{quote}

\begin{cousinweb}
\textbf{\bn{দক্ষতা}} (\emph{dôkkhôtā}) is Sanskrit \emph{dakṣatā}, dexterity or ability (Monier-Williams), from \emph{dakṣa} 'able'.
\end{cousinweb}

\begin{grammarlens}[title={What you've built}]
One more word for saying what happens if.
\end{grammarlens}

\subsection*{Your turn}
\begin{itemize}
\raggedright
  \item \emph{Say it:} \emph{dôkkhôtā}
  \item \emph{Say it:} \emph{dôkkhôtā}, once more
  \item \emph{Say it:} \emph{dôkkhôtā thākle chākri pāoyā sôhôj.} — If you have the skills, it's easy to get a job.
  \item \emph{Recall:} say the Bengali for effective, then the Bengali for out of order, then say \emph{dôkkhôtā} again
\end{itemize}

\begin{tcolorbox}[breakable,colback=teal!4,colframe=teal!35!black,title={Before you move on}]
What does \bn{দক্ষতা} mean? (\textquotedblleft{}Skill; competence\textquotedblright{}.) Say it once more.
\end{tcolorbox}
\section[\texorpdfstring{\bn{অনুশীলন} \bn{করা} (ônushīlôn kôrā)}{অনুশীলন করা (ônushīlôn kôrā)}]{\bn{অনুশীলন} \bn{করা} (ônushīlôn kôrā) — to practise}

\label{lesson:BN-C381-onushilon-kora}



\begin{quote}
Before the new one: say the Bengali for out of order, then the Bengali for skill.
\end{quote}

\subsection*{You'll want to know: \bn{অনুশীলন} \bn{করা}}
\textbf{\bn{অনুশীলন} \bn{করা}} — \emph{ônushīlôn kôrā} — \textquotedblleft{}to practise\textquotedblright{}.

\begin{quote}
\bn{রোজ} \bn{অনুশীলন} \bn{করলে} \bn{তুমি} \bn{ভালো} \bn{খেলবে।}
\end{quote}

\begin{quote}
\emph{roj ônushīlôn kôrle tumi bhālo khelbe.}
\end{quote}

\begin{quote}
If you practise every day, you'll play well.
\end{quote}

Past (he/she …): \bn{তিনি} \bn{সারাদিন} \bn{অনুশীলন} \bn{করলেন}, \emph{tini sārādin ônushīlôn kôrlen}.

\begin{cousinweb}
\textbf{\bn{অনুশীলন} \bn{করা}} (\emph{ônushīlôn kôrā}) is Sanskrit \emph{anuśīlana}, constant practice or study (Monier-Williams), with \bn{করা}: to do something again and again to get it right.
\end{cousinweb}

\begin{grammarlens}[title={What you've built}]
One more word for saying what happens if.
\end{grammarlens}

\subsection*{Your turn}
\begin{itemize}
\raggedright
  \item \emph{Say it:} \emph{ônushīlôn kôrā}
  \item \emph{Say it:} \emph{ônushīlôn kôrā}, once more
  \item \emph{Say it:} \emph{roj ônushīlôn kôrle tumi bhālo khelbe.} — If you practise every day, you'll play well.
  \item \emph{Recall:} say the Bengali for out of order, then the Bengali for skill, then say \emph{ônushīlôn kôrā} again
\end{itemize}

\begin{tcolorbox}[breakable,colback=teal!4,colframe=teal!35!black,title={Before you move on}]
What does \bn{অনুশীলন} \bn{করা} mean? (\textquotedblleft{}To practise\textquotedblright{}.) Say it once more.
\end{tcolorbox}
